\] for sufficiently large K at fixed m. This is impossible, so \(P(\lambda)=0\). For \(m=4\), a sufficient condition is \(\lambda^{K/4}(1-\lambda^{-K})^2>64\). It holds for \(K\geq11\): \(\lambda>5\), \(5^{11/4}>80\), and \(1-5^{-11}>9/10\) give a strict lower bound of 64.8. A nonzero polynomial of coefficient \(\ell^1\)-norm at most four cannot vanish at \(\lambda>3\), since its leading term strictly dominates the sum of all lower terms. Thus in this case the polynomial was zero. End of proof.

It follows that the \(2K(K-1)\) vectors \[
 \epsilon v_t+\delta v_s,\qquad 0\leq t<s<K,\quad
 \epsilon,\delta\in\{-1,1\},                                  \tag{FN.11}
\] are nonzero and distinct for \(K\geq11\). Any collision would give a balanced four-letter relation, forcing the original unordered signed pair. The matrices \(W(v)\) form an orthonormal basis for the full normalized Hilbert--Schmidt inner product, also for composite odd \(N\), because \(\operatorname{Tr}W(v)=N\mathbf1_{v=0}\).

Each of the four Weyl coefficients in \([B_t,B_s]\) has modulus \(a_{s-t}\). Different unordered time pairs have disjoint support by (FN.11). Since \(c_t-c_s=(s-t)/K\), exact orthogonality gives \[
 \|[G,\overline B]\|_{2,N}^2
 =\frac4{K^4}\sum_{j=1}^{K-1}(K-j)j^2a_j^2
 =\frac2{K^3}\sum_{j=1}^{K-1}j(K-j)a_j^2
 \leq\frac9K.                                                  \tag{FN.12}
\] Here \(\|A\|_{2,N}^2=\operatorname{Tr}(A^*A)/N\). The middle identity uses \(a_{K-j}=a_j\); the last uses \(j(K-j)\leq K^2/4\), \(a_j^2\leq2a_j\), and (FN.7).

The operator norm needs no orthogonality. Four Weyl coefficients, (FN.7), and (FN.8) give \[
 \|[G,B]\|_{\rm op}<\frac{26}{K},\qquad
 \|[B,\overline B]\|_{\rm op}<\frac{36}{K}.                     \tag{FN.13}
\] For a parity-preserving A, positivity of A*A implies \[
 \|A\|_{2,-}\leq\sqrt{\frac{2N}{N-1}}\,\|A\|_{2,N}.             \tag{FN.14}
\] No orthogonality of Weyl operators restricted to the odd subspace is assumed.

\subsection{A finite normal form without a small generator}\label{floquet:a-finite-normal-form-without-a-small-generator}

For finite Hermitian A,D, real \(\varepsilon\), and any unitarily invariant norm, \[
 \|e^{i\varepsilon A}e^{i\varepsilon D}
           -e^{i\varepsilon(A+D)}\|
 \leq\frac{\varepsilon^2}{2}\|[A,D]\|.                          \tag{FN.15}
\] This is the usual two-factor commutator estimate, not a new general product formula. Differentiate \(e^{i(\varepsilon-t)(A+D)}e^{itA}e^{itD}\). After factoring out unitaries, the integrand is \(e^{itA}De^{-itA}-D\). Integrate its derivative once more; the triangular integration domain has area \(\varepsilon^2/2\). Negative \(\varepsilon\) follows by replacing \(A,D\) with \(-A,-D\). In particular, \(|\varepsilon|\|A\|\) need not be small.

Set \(V_\varepsilon=e^{-i\varepsilon G}\). By (FN.4), \[
 V_\varepsilon U_\varepsilon V_\varepsilon^*
 =e^{-i\varepsilon G}e^{i\varepsilon B}
       e^{i\varepsilon C(G)}U_0.
\] Apply (FN.15) first to \(-G,B\), then to \(B-G,C(G)\). Their total sum is \(\bar B\), and \([B-G,C(G)]=[B-G,\bar B]\). Equations (FN.12)--(FN.14) give \[
 \begin{split}
 \|V_\varepsilon U_\varepsilon V_\varepsilon^*
       -e^{i\varepsilon\overline B}U_0\|_{2,-}
 &\leq\frac{\varepsilon^2}{2}
   \left(\frac{62}{K}+3\sqrt{\frac{2N}{(N-1)K}}\right)\\
 &\leq\frac{12\varepsilon^2}{\sqrt K},\qquad K\geq11.
 \end{split}\tag{FN.16}
\] For the last step it even suffices to use \(N\geq7\) and \(31/\sqrt{11}+(3/2)\sqrt{7/3}<12\). This proves (FN.D) for all real \(\varepsilon\). G may have operator norm of order K. Expanding its exponential in powers of \(\varepsilon\) and requiring \(\varepsilon K\) small would incorrectly discard the critical regime. Estimate (FN.15) avoids that requirement.

Since \(\bar B\) commutes with \(U_{0}\) and \(U_0^K=I\), telescoping unitary powers in the same norm yields, for any integer m, \[
 \left|\tau_-(U_\varepsilon^{mK})
       -\tau_-(e^{im\varepsilon K\overline B})\right|
 \leq12|m|\varepsilon^2\sqrt K,\qquad
 \tau_-(A)=d_N^{-1}\operatorname{Tr}_{-}A.                      \tag{FN.17}
\] For negative m use adjoints. Cauchy--Schwarz for the normalized trace justifies the passage from the norm to traces.

\subsection{A self-contained global law for the averaged observable}\label{floquet:a-self-contained-global-law-for-the-averaged-observable}

The remaining input is established here directly on the whole odd subspace; no unproved claim about the old arithmetic-moment corpus is imported. Define \[
 X_K=\sqrt{K/2}\,\overline B=(2K)^{-1/2}\sum_{j=0}^{K-1}B_j.
\]

\textbf{Lemma 2 (Gaussian average).} For each fixed integer \(r\geq1\), \[
 \tau_-(X_K^{2r})=(2r-1)!!+O_r(K^{-1})+O_r(K^r/N),               \tag{FN.18}
\] and for each fixed odd m, eventually \[
 |\tau_-(X_K^m)|\leq (2K)^{m/2}/(N-1).                         \tag{FN.19}
\] The empirical measures of \(X_{K}\) converge weakly and in quadratic Wasserstein distance to the standard normal law.

\textbf{Proof.} In the full Weyl representation \(\operatorname{Tr}(RW(v))=1\) for every v: the unique fixed point \(x=-b/2\) of RW(a,b) contributes phase 1. Thus the normalized odd trace of a word \(e_{N}\)(\(\phi\))W(v) is \[
 e_N(\phi)\frac{N\mathbf1_{v=0}-1}{N-1}.                       \tag{FN.20}
\] Expanding \(X_K^m\), the noncentral part is bounded by \((2K)^{m/2}/(N-1)\). The central part is the weighted count of the reciprocal relations (FN.9), normalized by \([N/(N-1)](2K)^{-m/2}\).

Here is the count including structurally forced nonbalanced relations. For a polynomial of coefficient \(\ell^1\)-norm at most m vanishing at \(\lambda\), a gap \(d>2\log_\lambda m\) separates two independently vanishing pieces. Write \(P=P_{\rm low}+Z^{a+d}P_{\rm high}\), with \(\deg P_{\rm low}\leq a\) and nonnegative exponents in \(P_{\rm high}\). If \(P_{\rm high}(\lambda)\ne0\), its algebraic-integer norm is a nonzero integer. Since \(|P_{\rm high}(\lambda^{-1})|\leq m\), we have \(|P_{\rm high}(\lambda)|\geq1/m\). Cancellation with \(|P_{\rm low}(\lambda)|\leq m\lambda^a\) would force \(\lambda^d\leq m^2\), a contradiction.

Repeatedly splitting gives components with span bounded in terms of m, for example by \(D_m=(m-1)(\lceil2\log_\lambda m\rceil+1)\). Every nonzero component has at least three signed occurrences: two distinct positive powers cannot cancel with signs of magnitude one. Restore previously cancelled opposite pairs as two-occurrence components. A nonbalanced 2r-letter word therefore has at most r-1 components. For each of finitely many partitions of the 2r positions, signs and bounded relative offsets, each component has at most K translations on the cyclic time set. There are only \(O_r(K^{r-1})\) such words. There is no additional factor K from choosing a cut: the count is over the component translations themselves, and the largest-gap cut was used only to prove bounded width. Their normalized contribution is \(O_r(1/K)\).

For odd m, a lifted P is an integer multiple of \(Z^2-6Z+1\), so \(P(1)\) is divisible by -4. But \(P(1)\) is the sum of an odd number of signs. Thus no central word exists once Lemma 1 applies, proving (FN.19).

A balanced \(2r\)-word pairs identical time labels with opposite signs. Assignments with repeated pair labels cause only \(O_r(K^{r-1})\) overcount. Each matching of the \(2r\) positions has \(K^r\) pair-label choices and \(2^r\) orientations. Its Weyl phase is a product, over crossing pairs, of factors \(e_N(q(h^a-h^{-a}))\) or their conjugates. Summing the signed symplectic products in word order shows that disjoint or nested pairs cancel, while a crossing pair retains one full symplectic product.

Each crossing difference of independent uniform pair labels is uniform modulo \(K\). Independence of the different crossing factors is neither true in general nor needed. The inequality \(|\prod_l u_l-1|\leq\sum_l|u_l-1|\), for \(|u_l|=1\), and (FN.7) give average phase error at most nine times the number of crossings divided by \(K\). Hence each matching contributes \(1+O_r(1/K)\) after normalization. There are \((2r-1)!!\) matchings. Together with the nonbalanced count and (FN.20), this proves (FN.18).

